% \subsection{Velocidades linear e angular}
% \label{sec_dinAcoplada_veloc_linear_angular}

% As velocidades linear ${}^{B}\!\vec v_{G_i}$ e angular ${}^{B}\!\vec\omega_i$
% do elo $i$ são obtidas a partir do jacobiano geométrico avaliado no respectivo
% centro de massa \cite{craig2005,Wang2003KinematicsDynamics}.
% Escrevendo o jacobiano na forma espacial \cite{murray1994mathematical},
% tem-se um operador $J_{\text{sp},i}(\vec y\, \, )$ que
% relaciona diretamente as velocidades generalizadas
% $\dot{\vec y\, \, }$ com as velocidades cinemáticas empregadas na energia
% cinética:

% \begin{equation}
% \begin{bmatrix}
% {}^{B}\!\vec v_{G_i} \\
% {}^{B}\!\vec\omega_i
% \end{bmatrix}
% =
% J_{\text{sp},i}(\vec y\,)\,\dot{\vec y\,}.
% \label{eq_jacobianos_vGi1}
% \end{equation}

% Em seguida, decompõe-se o jacobiano espacial em parte linear e parte angular

% \begin{equation}
% J_{\text{sp},i}(\vec y\,)
% =
% \begin{bmatrix}
% J_{v,i}(\vec y\,) \\
% J_{\omega,i}(\vec y\,)
% \end{bmatrix},
% \label{eq_jacobianos_vGi2}
% \end{equation}

% Em que o termo $J_{v,i}(\vec y\, )$ é a parte linear, enquanto
% $J_{\omega,i}(\vec y\, )$ é a parte angular do jacobiano espacial do elo $i$.

% Como discutido em \cite{sciavicco_siciliano2010}, no caso de juntas
% revolutas as colunas de $J_{v,i}$ e $J_{\omega,i}$ associadas às coordenadas
% articulares podem ser escritas a partir dos eixos de junta e das posições
% dos centros de massa.
% %%%%%%%%%%%%%%%%%%%%%%%%%%%%%%%%%%%%%%%%%%%%%%%%

% Se ${}^{B}\!\vec z_j$ denota o eixo de rotação da junta $j$ expresso em $\{B\}$,
% ${}^{B}\!\vec r_{O_j}$ a posição da origem da junta $j$ em $\{B\}$ e
% ${}^{B}\!\vec r_{G_i}$ a posição do centro de massa do elo $i$ em $\{B\}$,
% então ${}^{B}\!\vec v_{G_i}$ representa a velocidade linear no \gls{cdm} do elo $i$,
% e ${}^{B}\!\vec\omega_i$ representa a velocidade angular do elo $i$, ambas expressas em $\{B\}$.
% Para uma junta revoluta $j$ que atua sobre o elo $i$, a coluna correspondente do
% jacobiano espacial avaliado no seu \gls{cdm} $G_i$ é dada por

% \begin{equation}
% J_{v,i}^{(j)}(\vec y\, )
% =
% {}^{B}\!\vec z_j \times
% \bigl({}^{B}\!\vec r_{G_i}(\vec y\, ) - {}^{B}\!\vec r_{O_j}(\vec y\, )\bigr),
% \qquad
% J_{\omega,i}^{(j)}(\vec y\, )
% =
% {}^{B}\!\vec z_j,
% \label{eq_colunas_J_sp_revoluta}
% \end{equation}

% para $j\le i$; para $j>i$, essas colunas são nulas, pois a junta $j$ não influencia o elo $i$.
% Para evidenciar a estrutura do jacobiano espacial $J_{\text{sp},i}$,
% escrevem-se explicitamente as colunas não nulas para $i=1,2,3$.
% Para $i=1$:

% \begin{equation}
% J_{\text{sp},1}(\vec y\, )
% =
% \begin{bmatrix}
% {}^{B}\!\vec z_1 \times \bigl({}^{B}\!\vec r_{G_1}-{}^{B}\!\vec r_{O_1}\bigr) & \vec 0 & \vec 0\\[2pt]
% {}^{B}\!\vec z_1 & \vec 0 & \vec 0
% \end{bmatrix}.
% \label{eq_Jsp_i1}
% \end{equation}

% Para $i=2$:
% \begin{equation}
% J_{\text{sp},2}(\vec y\, )
% =
% \begin{bmatrix}
% {}^{B}\!\vec z_1 \times \bigl({}^{B}\!\vec r_{G_2}-{}^{B}\!\vec r_{O_1}\bigr) &
% {}^{B}\!\vec z_2 \times \bigl({}^{B}\!\vec r_{G_2}-{}^{B}\!\vec r_{O_2}\bigr) &
% \vec 0\\[2pt]
% {}^{B}\!\vec z_1 &
% {}^{B}\!\vec z_2 &
% \vec 0
% \end{bmatrix}.
% \label{eq_Jsp_i2}
% \end{equation}

% Para $i=3$:
% \begin{equation}
% J_{\text{sp},3}(\vec y\, )
% =
% \begin{bmatrix}
% {}^{B}\!\vec z_1 \times \bigl({}^{B}\!\vec r_{G_3}-{}^{B}\!\vec r_{O_1}\bigr) &
% {}^{B}\!\vec z_2 \times \bigl({}^{B}\!\vec r_{G_3}-{}^{B}\!\vec r_{O_2}\bigr) &
% {}^{B}\!\vec z_3 \times \bigl({}^{B}\!\vec r_{G_3}-{}^{B}\!\vec r_{O_3}\bigr)\\[2pt]
% {}^{B}\!\vec z_1 &
% {}^{B}\!\vec z_2 &
% {}^{B}\!\vec z_3
% \end{bmatrix}.
% \label{eq_Jsp_i3}
% \end{equation}

% Os termos dependentes de velocidade (Coriolis e centrífugos) decorrem da aplicação
% das equações de Lagrange à energia cinética do sistema e são reunidos no termo
% $C(\cdot)\,\dot{\cdot}$ nas equações de movimento \cite{Geradin2004,Wilde2018}.
% Assim, o foco recai sobre a construção de $M(\cdot)$ por meio dos jacobianos
% e dos operadores de inércia espacial.

% Em ambiente de microgravidade, o centro de massa do sistema praticamente coincide com o centro de gravidade e o torque gravitacional interno em torno do centro de massa é desprezível \cite{wie2008space}.
% Além disso, para uma órbita circular, a contribuição da energia potencial gravitacional na dinâmica de atitude é praticamente constante.
% Assim, no modelo adotado, toma-se a energia potencial generalizada como
% $V(\vec y\, ) = 0$ e a Lagrangiana reduz-se
% a $L(\vec y\, ,\dot{\vec y\, }) = T(\vec y\, ,\dot{\vec y\, })$.
% O conjunto espaçonave--\gls{mr} é, então, descrito diretamente pelas equações de Lagrange em sua forma padrão \cite{murray1994mathematical}:

% \begin{equation}
% M(\vec y\, )\,\ddot{\vec y\, }
% +
% C(\vec y\, ,\dot{\vec y\, })\,\dot{\vec y\, }
% =
% \vec\tau,
% \label{eq_lagrange_formato_padrao}
% \end{equation}

% Em que $\vec y\, $ é o vetor de coordenadas generalizadas, $M(\vec y\, )$ é a
% matriz de inércia generalizada e $C(\vec y\, ,\dot{\vec y\, })\,\dot{\vec y\, }$
% agrega os termos de Coriolis e centrífugos.
% O foco desta seção é a construção de $M(\vec y\, )$ a partir das
% contribuições de massa e inércia da espaçonave e dos elos do manipulador.
% %%%%%%%%%%%%%%%%%%%

% % As coordenadas generalizadas são organizadas de modo a reunir, em um único
% % vetor, os parâmetros de atitude da espaçonave e as coordenadas articulares do manipulador:

% % \begin{equation}
% % \vec y\, \in \mathbb{R}^{3+n}
% % =
% % \begin{bmatrix}
% % \vec\eta_B \\[2pt]
% % \vec\xi\end{bmatrix},
% % \qquad
% % \dot{\vec y\, } \in \mathbb{R}^{3+n}
% % =
% % \begin{bmatrix}
% % \dot{\vec\eta}_B \\[2pt]
% % \dot{\vec\xi}
% % \end{bmatrix},
% % \label{eq_q_generalizado_espaconave_mr}
% % \end{equation}

% % Em que $\vec\eta_B \in \mathbb{R}^{3}$ representa os ângulos de Euler da
% % espaçonave e $\vec\xi \in \mathbb{R}^{n}$ contém as coordenadas articulares
% % do manipulador (neste trabalho, $n=3$).

% As coordenadas generalizadas são organizadas de modo a reunir, em um único
% vetor, os parâmetros de atitude da espaçonave e as coordenadas articulares do
% manipulador:

% \begin{equation}
% \vec y\, \in \mathbb{R}^{3+n}
% =
% \begin{bmatrix}
% \vec\eta_B \\[2pt]
% \vec\xi\end{bmatrix},
% \qquad
% \dot{\vec y\, } \in \mathbb{R}^{3+n}
% =
% \begin{bmatrix}
% \dot{\vec\eta}_B \\[2pt]
% \dot{\vec\xi}
% \end{bmatrix},
% \label{eq_q_generalizado_espaconave_mr}
% \end{equation}

% Em que $\vec\eta_B = [\phi_1, \phi_2, \phi_3]^T$ representa os ângulos de Euler da
% espaçonave (sequência 3--2--1) e $\vec\xi \in \mathbb{R}^{n}$ contém as coordenadas articulares do manipulador.

% %%%%%%%%%%%%
% % A relação cinemática entre as taxas dos ângulos de Euler e a velocidade
% % angular do corpo é dada pelas equações cinemáticas diferenciais
% % \cite{Schaub2010AttitudeDynamics}, que, para a sequência 3--2--1,
% % podem ser escritas como

% % \begin{equation}
% % \dot{\vec\eta}_B
% % =
% % B(\vec\eta_B)\,{}^{B}\!\vec\omega_B,
% % \qquad
% % \vec\eta_B=
% % \begin{bmatrix}
% % \phi\\ \xi\\ \psi
% % \end{bmatrix},
% % \label{eq_euler321_kinematics}
% % \end{equation}

% % Em que $\phi$, $\xi$ e $\psi$ são, respectivamente, os ângulos de rolagem,
% % arfagem e guinada.
% A relação cinemática entre as taxas dos ângulos de Euler e a velocidade
% angular do corpo, ${}^{B}\!\vec\omega_B$, é dada diretamente pela matriz de transformação $\mathbf{S}(\vec\eta_B)$, conforme apresentado na Equação~\eqref{eq_cinematica_euler}:

% \begin{equation}
% {}^{B}\!\vec\omega_B
% =
% \mathbf{S}(\vec\eta_B)\,\dot{\vec\eta}_B,
% \label{eq:euler_direct_kinematics}
% \end{equation}

% Em que a matriz $\mathbf{S}(\vec\eta_B)$ é definida como \cite{wie2008space}:
% \begin{equation}
% \mathbf{S}(\vec\eta_B)
% =
% \begin{bmatrix}
% 1 & 0 & -\sin\phi_2 \\
% 0 & \cos\phi_1 & \sin\phi_1 \cos\phi_2 \\
% 0 & -\sin\phi_1 & \cos\phi_1 \cos\phi_2
% \end{bmatrix}.
% \label{eq:S_matrix_definition}
% \end{equation}
% %%%%%%%%%%%%%%%%%%%%%%%%%%%
% Para a sequência 3--2--1, a matriz $B(\vec\eta_B)\in\mathbb{R}^{3\times 3}$ é

% \begin{equation}
% B(\vec\eta_B)
% =
% \begin{bmatrix}
% 1 & \sin\phi\,\tan\xi & \cos\phi\,\tan\xi\\
% 0 & \cos\phi             & -\sin\phi\\
% 0 & \sin\phi\,\sec\xi & \cos\phi\,\sec\xi
% \end{bmatrix}.
% \label{eq_B_euler321}
% \end{equation}

% Neste trabalho, utiliza-se a forma inversa, de modo a expressar a velocidade
% angular da espaçonave em função das velocidades generalizadas de atitude:

% \begin{equation}
% {}^{B}\!\vec\omega_B
% =
% \mathbf{T}_\omega(\vec\eta_B)\,\dot{\vec\eta}_B,
% \qquad
% \mathbf{T}_\omega(\vec\eta_B)=B^{-1}(\vec\eta_B),
% \label{eq_euler321_inverse_kinematics}
% \end{equation}

% com
% \begin{equation}
% \mathbf{T}_\omega(\vec\eta_B)
% =
% \begin{bmatrix}
% 1 & 0        & -\sin\xi\\
% 0 & \cos\phi & \sin\phi\,\cos\xi\\
% 0 & -\sin\phi& \cos\phi\,\cos\xi
% \end{bmatrix}.
% \label{eq_Tomega_euler321}
% \end{equation}


% %%%%%%%%%%%%%%%%
% % Ao reunir a velocidade angular da espaçonave e as velocidades articulares em um único vetor, define-se

% % \begin{equation}
% % \dot{\vec s}
% % =
% % \begin{bmatrix}
% % {}^{B}\!\vec\omega_B \\[2pt]
% % \dot{\vec\xi}
% % \end{bmatrix}
% % \in \mathbb{R}^{3+n}.
% % \label{eq_def_s_dot}
% % \end{equation}

% Ao reunir a velocidade angular da espaçonave e as velocidades articulares em um único vetor de quase-velocidades, define-se:

% \begin{equation}
% \dot{\vec s}
% =
% \begin{bmatrix}
% {}^{B}\!\vec\omega_B \\[2pt]
% \dot{\vec\xi}
% \end{bmatrix}
% \in \mathbb{R}^{3+n}.
% \label{eq:def_s_dot}
% \end{equation}
% %%%%%%%%%%%%%%%%%%%%%%%%%

% % Como ${}^{B}\!\vec\omega_B=\mathbf{T}_\omega(\vec\eta_B)\,\dot{\vec\eta}_B$ e $\dot{\vec\xi}=I_n\,\dot{\vec\xi}$, obtém-se

% % \begin{equation}
% % \dot{\vec s}
% % =
% % \underbrace{
% % \begin{bmatrix}
% % \mathbf{T}_\omega(\vec\eta_B) & 0 \\[2pt]
% % 0 & I_{n}
% % \end{bmatrix}}_{\displaystyle \mathbf{J}_s(\vec\eta_B)}
% % \dot{\vec y\, },
% % \qquad
% % \dot{\vec y\, }
% % =
% % \begin{bmatrix}
% % \dot{\vec\eta}_B \\[2pt]
% % \dot{\vec\xi}
% % \end{bmatrix}.
% % \label{eq_Js_mapping}
% % \end{equation}

% Substituindo \eqref{eq:euler_direct_kinematics} em \eqref{eq:def_s_dot} e considerando que $\dot{\vec\xi}=I_n\,\dot{\vec\xi}$, obtém-se o mapeamento entre as velocidades generalizadas $\dot{\vec y\,}$ e o vetor cinemático $\dot{\vec s}$:

% \begin{equation}
% \dot{\vec s}
% =
% \underbrace{
% \begin{bmatrix}
% \mathbf{S}(\vec\eta_B) & 0 \\[2pt]
% 0 & I_{n}
% \end{bmatrix}}_{\displaystyle \mathbf{J}_s(\vec\eta_B)}
% \underbrace{
% \begin{bmatrix}
% \dot{\vec\eta}_B \\[2pt]
% \dot{\vec\xi}
% \end{bmatrix}}_{\displaystyle \dot{\vec y\,}}.
% \label{eq:Js_mapping}
% \end{equation}

% A matriz $\mathbf{J}_s(\vec\eta_B)$ é o Jacobiano de transformação que relaciona as derivadas das coordenadas generalizadas com as velocidades físicas (angular da espaçonave e das juntas).
% Com essa definição, a energia cinética total do conjunto pode ser escrita como

% \begin{equation}
% T =
% \frac{1}{2}\,\dot{\vec s}^{\,T}\,M_s(\vec y\, )\,\dot{\vec s},
% \label{eq_T_matriz_s}
% \end{equation}

% Em que $M_s(\vec y\, )$ é a matriz de inércia generalizada escrita no espaço do vetor
% cinemático $\dot{\vec s}$.

% %%%%%%%%%%%%%%%%%%%%%%%%%%%%%%%


% Para cada elo, a matriz $M_{i,\text{spt}}$ é um operador de inércia
% espacial $6\times 6$ que combina, em um único bloco, os efeitos
% translacionais e rotacionais em torno do centro de massa expresso em
% $\{B\}$.
% Se ${}^{B}\!p_{G_i}$ denota a posição do centro de massa do elo em relação à
% origem de $\{B\}$ e ${}^{B}\!I_i$ o matriz de inércia do elo em torno desse
% centro, a forma padrão do operador de inércia espacial é

% \begin{equation}
% M_{i,\text{spt}}
% =
% \begin{bmatrix}
% m_i I_3 & -m_i\,S({}^{B}\!p_{G_i})\\[2pt]
% m_i\,S({}^{B}\!p_{G_i}) &
% {}^{B}\!I_i - m_i\,S({}^{B}\!p_{G_i})\,S({}^{B}\!p_{G_i})^{T}
% \end{bmatrix},
% \end{equation}

% Em que $m_i$ é a massa do $i$-ésimo elo, $I_3$ é a matriz identidade $3\times 3$,
% ${}^{B}\!p_{G_i}$ é o vetor posição do centro de massa do elo $i$
% em relação à origem do referencial $\{B\}$, ${}^{B}\!I_i$
% é o matriz de inércia do elo $i$ em torno do seu centro de massa, expresso em
% $\{B\}$, e $S(\cdot)$ é o operador matricial antissimétrico associado ao produto
% vetorial, tal que $S(\vec a)\,\vec b = \vec a \times \vec b$ para quaisquer
% $\vec a,\vec b$.

% De maneira análoga, a espaçonave é descrita por um operador de inércia
% espacial $M_{\text{espaç},\text{spt}}$.
% No formalismo de velocidades espaciais \cite{featherstone2008}, uma matriz de
% inércia espacial relaciona a velocidade espacial
% $\vec v_{\text{sp}} = [\,{}^{B}\!\vec v^{\,T}\;{}^{B}\!\vec\omega^{\,T}\,]^{T}$
% com a quantidade de movimento espacial
% $\vec h_{\text{sp}}$ por meio de
% $\vec h_{\text{sp}} = M_{\text{spt}}\,\vec v_{\text{sp}}$.

% Para a espaçonave, $M_{\text{espaç},\text{spt}}$ é construído a partir da massa
% $m_B$ e do matriz ${}^{B}\!I_B$ em torno do seu centro de massa, em completa
% analogia com $M_{i,\text{spt}}$ definido para os elos.

% A energia cinética total do conjunto espaçonave--\gls{mr} resulta da soma das
% energias cinéticas de cada elo e da espaçonave.
% Utilizando a forma espacial para cada elo, tem-se

% \begin{equation}
% T
% =
% \sum_{i=1}^{n}
% \frac{1}{2}
% \bigl(J_{\text{sp},i}(\vec y\, )\,\dot{\vec y\, }\bigr)^{T}
% M_{i,\text{spt}}
% \bigl(J_{\text{sp},i}(\vec y\, )\,\dot{\vec y\, }\bigr)
% +
% \frac{1}{2}\,
% \vec v_{\text{espaç,sp}}^{\,T}\,
% M_{\text{espaç},\text{spt}}\,
% \vec v_{\text{espaç,sp}},
% \end{equation}

% Em que $\vec v_{\text{espaç,sp}}$ é a velocidade espacial da espaçonave, escrita
% também como combinação linear das velocidades generalizadas $\dot{\vec y\, }$.
% Comparando essa expressão com a forma quadrática de \eqref{eq_T_quadratica_Mq1}, obtém-se

% \begin{equation}
% T =
% \frac{1}{2}\,\dot{\vec y\, }^{\,T}\,M(\vec y\, )\,\dot{\vec y\, },
% \end{equation}

% de modo que a matriz de inércia generalizada $M_s(\vec y\, )$ é dada pela soma das
% contribuições espaciais de todos os elos, ponderadas pelos respectivos
% jacobianos espaciais, acrescida da contribuição inercial da espaçonave:

% \begin{equation}
% M_s(\vec y\, ) =
% \sum_{i=1}^{n}
% J_{\text{sp},i}^{T}(\vec y\, )\,
% M_{i,\text{spt}}\,
% J_{\text{sp},i}(\vec y\, )
% +
% M_{\text{espaç},\text{spt}}.
% \label{eq_Ms_total_formula}
% \end{equation}


% Cada termo
% $J_{\text{sp},i}^{T}(\vec y\, \, \, )\,M_{i,\text{spt}}\,J_{\text{sp},i}(\vec y\, \, \, )$
% é, portanto, a contribuição espacial do elo $i$ para a matriz de inércia
% generalizada $M_s(\vec y\, )$.
% Nessa expressão, cada jacobiano espacial $J_{\text{sp},i}(\vec y\, \, )$ mapeia as velocidades generalizadas $\dot{\vec y\, \, }$ na velocidade espacial (linear e angular) do centro de massa do elo $i$ em $\{B\}$.


\section{Velocidades linear e angular}
\label{sec_dinAcoplada_veloc_linear_angular}

As velocidades linear ${}^{B}\!\vec v_{G_i}$ e angular ${}^{B}\!\vec\omega_i$
do elo $i$ são obtidas a partir do jacobiano geométrico avaliado no respectivo
centro de massa \cite{craig2005,Wang2003KinematicsDynamics}.
Escrevendo o jacobiano na forma espacial \cite{murray1994mathematical},
tem-se um operador $J_{\text{sp},i}(\vec y\, \, )$ que
relaciona diretamente as velocidades generalizadas
$\dot{\vec y\, \, }$ com as velocidades cinemáticas empregadas na energia
cinética:

\begin{equation}
\begin{bmatrix}
{}^{B}\!\vec v_{G_i} \\
{}^{B}\!\vec\omega_i
\end{bmatrix}
=
J_{\text{sp},i}(\vec y\,)\,\dot{\vec y\,}.
\label{eq_jacobianos_vGi1}
\end{equation}

Em seguida, decompõe-se o jacobiano espacial em parte linear e parte angular:

\begin{equation}
J_{\text{sp},i}(\vec y\,)
=
\begin{bmatrix}
J_{v,i}(\vec y\,) \\
J_{\omega,i}(\vec y\,)
\end{bmatrix}.
\label{eq_jacobianos_vGi2}
\end{equation}

Em que o termo $J_{v,i}(\vec y\, )$ é a parte linear, enquanto
$J_{\omega,i}(\vec y\, )$ é a parte angular do jacobiano espacial do elo $i$.

Como discutido em \cite{sciavicco_siciliano2010}, no caso de juntas
revolutas as colunas de $J_{v,i}$ e $J_{\omega,i}$ associadas às coordenadas
articulares podem ser escritas a partir dos eixos de junta e das posições
dos centros de massa.
%%%%%%%%%%%%%%%%%%%%%%%%%%%%%%%%%%%%%%%%%%%%%%%%

Se ${}^{B}\!\vec z_j$ denota o eixo de rotação da junta $j$ expresso em $\{B\}$,
${}^{B}\!\vec r_{O_j}$ a posição da origem da junta $j$ em $\{B\}$ e
${}^{B}\!\vec r_{G_i}$ a posição do centro de massa do elo $i$ em $\{B\}$,
então ${}^{B}\!\vec v_{G_i}$ representa a velocidade linear no \gls{cdm} do elo $i$,
e ${}^{B}\!\vec\omega_i$ representa a velocidade angular do elo $i$, ambas expressas em $\{B\}$.
Para uma junta revoluta $j$ que atua sobre o elo $i$, a coluna correspondente do
jacobiano espacial avaliado no seu \gls{cdm} $G_i$ é dada por

\begin{equation}
J_{v,i}^{(j)}(\vec y\, )
=
{}^{B}\!\vec z_j \times
\bigl({}^{B}\!\vec r_{G_i}(\vec y\, ) - {}^{B}\!\vec r_{O_j}(\vec y\, )\bigr),
\qquad
J_{\omega,i}^{(j)}(\vec y\, )
=
{}^{B}\!\vec z_j,
\label{eq_colunas_J_sp_revoluta}
\end{equation}

para $j\le i$; para $j>i$, essas colunas são nulas, pois a junta $j$ não influencia o elo $i$.
Para evidenciar a estrutura do jacobiano espacial $J_{\text{sp},i}$,
escrevem-se explicitamente as colunas não nulas para $i=1,2,3$.
Para $i=1$:

\begin{equation}
J_{\text{sp},1}(\vec y\, )
=
\begin{bmatrix}
{}^{B}\!\vec z_1 \times \bigl({}^{B}\!\vec r_{G_1}-{}^{B}\!\vec r_{O_1}\bigr) & \vec 0 & \vec 0\\[2pt]
{}^{B}\!\vec z_1 & \vec 0 & \vec 0
\end{bmatrix}.
\label{eq_Jsp_i1}
\end{equation}

Para $i=2$:
\begin{equation}
J_{\text{sp},2}(\vec y\, )
=
\begin{bmatrix}
{}^{B}\!\vec z_1 \times \bigl({}^{B}\!\vec r_{G_2}-{}^{B}\!\vec r_{O_1}\bigr) &
{}^{B}\!\vec z_2 \times \bigl({}^{B}\!\vec r_{G_2}-{}^{B}\!\vec r_{O_2}\bigr) &
\vec 0\\[2pt]
{}^{B}\!\vec z_1 &
{}^{B}\!\vec z_2 &
\vec 0
\end{bmatrix}.
\label{eq_Jsp_i2}
\end{equation}

Para $i=3$:
\begin{equation}
J_{\text{sp},3}(\vec y\, )
=
\begin{bmatrix}
{}^{B}\!\vec z_1 \times \bigl({}^{B}\!\vec r_{G_3}-{}^{B}\!\vec r_{O_1}\bigr) &
{}^{B}\!\vec z_2 \times \bigl({}^{B}\!\vec r_{G_3}-{}^{B}\!\vec r_{O_2}\bigr) &
{}^{B}\!\vec z_3 \times \bigl({}^{B}\!\vec r_{G_3}-{}^{B}\!\vec r_{O_3}\bigr)\\[2pt]
{}^{B}\!\vec z_1 &
{}^{B}\!\vec z_2 &
{}^{B}\!\vec z_3
\end{bmatrix}.
\label{eq_Jsp_i3}
\end{equation}

Os termos dependentes de velocidade (Coriolis e centrífugos) decorrem da aplicação
das equações de Lagrange à energia cinética do sistema e são reunidos no termo
$C(\cdot)\,\dot{\cdot}$ nas equações de movimento \cite{Geradin2004,Wilde2018}.
Assim, o foco recai sobre a construção de $M(\cdot)$ por meio dos jacobianos
e dos operadores de inércia espacial.

Em ambiente de microgravidade, o centro de massa do sistema praticamente coincide com o centro de gravidade e o torque gravitacional interno em torno do centro de massa é desprezível \cite{wie2008space}.
Além disso, para uma órbita circular, a contribuição da energia potencial gravitacional na dinâmica de atitude é praticamente constante.
Assim, no modelo adotado, toma-se a energia potencial generalizada como
$V(\vec y\, ) = 0$ e a Lagrangiana reduz-se
a $L(\vec y\, ,\dot{\vec y\, }) = T(\vec y\, ,\dot{\vec y\, })$.
O conjunto espaçonave--\gls{mr} é, então, descrito diretamente pelas equações de Lagrange em sua forma padrão \cite{murray1994mathematical}:

\begin{equation}
M(\vec y\, )\,\ddot{\vec y\, }
+
C(\vec y\, ,\dot{\vec y\, })\,\dot{\vec y\, }
=
\vec\tau.
\label{eq_lagrange_formato_padrao}
\end{equation}

Em que $\vec y\, $ é o vetor de coordenadas generalizadas, $M(\vec y\, )$ é a
matriz de inércia generalizada e $C(\vec y\, ,\dot{\vec y\, })\,\dot{\vec y\, }$
agrega os termos de Coriolis e centrífugos.
O foco desta seção é a construção de $M(\vec y\, )$ a partir das
contribuições de massa e inércia da espaçonave e dos elos do manipulador.
%%%%%%%%%%%%%%%%%%%

As coordenadas generalizadas são organizadas de modo a reunir, em um único
vetor, os parâmetros de atitude da espaçonave e as coordenadas articulares do
manipulador:

\begin{equation}
\vec y\, \in \mathbb{R}^{3+n}
=
\begin{bmatrix}
\vec\eta_B \\[2pt]
\vec\xi\end{bmatrix},
\qquad
\dot{\vec y\, } \in \mathbb{R}^{3+n}
=
\begin{bmatrix}
\dot{\vec\eta}_B \\[2pt]
\dot{\vec\xi}
\end{bmatrix}.
\label{eq_q_generalizado_espaconave_mr}
\end{equation}

Em que $\vec\eta_B = [\phi_1, \phi_2, \phi_3]^T$ representa os ângulos de Euler da
espaçonave (sequência 3--2--1, conforme definido na Seção~\ref{sec_angulos_euler}) e $\vec\xi \in \mathbb{R}^{n}$ contém as coordenadas articulares do manipulador.

%%%%%%%%%%%%
A relação cinemática entre as taxas dos ângulos de Euler e a velocidade
angular do corpo, ${}^{B}\!\vec\omega_B$, é dada diretamente pela matriz de transformação $\mathbf{S}(\vec\eta_B)$, conforme apresentado na Equação~\eqref{eq_cinematica_euler}:

\begin{equation}
{}^{B}\!\vec\omega_B
=
\mathbf{S}(\vec\eta_B)\,\dot{\vec\eta}_B.
\label{eq:euler_direct_kinematics}
\end{equation}

Em que a matriz $\mathbf{S}(\vec\eta_B)$ é definida como \cite{wie2008space}:

\begin{equation}
\mathbf{S}(\vec\eta_B)
=
\begin{bmatrix}
1 & 0 & -\sin\phi_2 \\
0 & \cos\phi_1 & \sin\phi_1 \cos\phi_2 \\
0 & -\sin\phi_1 & \cos\phi_1 \cos\phi_2
\end{bmatrix}.
\label{eq:S_matrix_definition}
\end{equation}

Ao reunir a velocidade angular da espaçonave e as velocidades articulares em um único vetor de quase-velocidades, define-se:

\begin{equation}
\dot{\vec s}
=
\begin{bmatrix}
{}^{B}\!\vec\omega_B \\[2pt]
\dot{\vec\xi}
\end{bmatrix}
\in \mathbb{R}^{3+n}.
\label{eq:def_s_dot}
\end{equation}

Substituindo \eqref{eq:euler_direct_kinematics} em \eqref{eq:def_s_dot} e considerando que $\dot{\vec\xi}=I_n\,\dot{\vec\xi}$, obtém-se o mapeamento entre as velocidades generalizadas $\dot{\vec y\,}$ e o vetor cinemático $\dot{\vec s}$:

\begin{equation}
\dot{\vec s}
=
\underbrace{
\begin{bmatrix}
\mathbf{S}(\vec\eta_B) & 0 \\[2pt]
0 & I_{n}
\end{bmatrix}}_{\displaystyle \mathbf{J}_s(\vec\eta_B)}
\underbrace{
\begin{bmatrix}
\dot{\vec\eta}_B \\[2pt]
\dot{\vec\xi}
\end{bmatrix}}_{\displaystyle \dot{\vec y\,}}.
\label{eq:Js_mapping}
\end{equation}

A matriz $\mathbf{J}_s(\vec\eta_B)$ é o Jacobiano de transformação que relaciona as derivadas das coordenadas generalizadas com as velocidades físicas (angular da espaçonave e das juntas).
Com essa definição, a energia cinética total do conjunto pode ser escrita como:

\begin{equation}
T =
\frac{1}{2}\,\dot{\vec s}^{\,T}\,M_s(\vec y\, )\,\dot{\vec s}.
\label{eq_T_matriz_s}
\end{equation}

Em que $M_s(\vec y\, )$ é a matriz de inércia generalizada escrita no espaço do vetor
cinemático $\dot{\vec s}$.

Para cada elo, a matriz $M_{i,\text{spt}}$ é um operador de inércia
espacial $6\times 6$ que combina, em um único bloco, os efeitos
translacionais e rotacionais em torno do centro de massa expresso em
$\{B\}$.
Se ${}^{B}\!p_{G_i}$ denota a posição do centro de massa do elo em relação à
origem de $\{B\}$ e ${}^{B}\!I_i$ o matriz de inércia do elo em torno desse
centro, a forma padrão do operador de inércia espacial é:

\begin{equation}
M_{i,\text{spt}}
=
\begin{bmatrix}
m_i I_3 & -m_i\,S({}^{B}\!p_{G_i})\\[2pt]
m_i\,S({}^{B}\!p_{G_i}) &
{}^{B}\!I_i - m_i\,S({}^{B}\!p_{G_i})\,S({}^{B}\!p_{G_i})^{T}
\end{bmatrix}.
\end{equation}

Em que $m_i$ é a massa do $i$-ésimo elo, $I_3$ é a matriz identidade $3\times 3$,
${}^{B}\!p_{G_i}$ é o vetor posição do centro de massa do elo $i$
em relação à origem do referencial $\{B\}$, ${}^{B}\!I_i$
é o matriz de inércia do elo $i$ em torno do seu centro de massa, expresso em
$\{B\}$, e $S(\cdot)$ é o operador matricial antissimétrico associado ao produto
vetorial, tal que $S(\vec a)\,\vec b = \vec a \times \vec b$ para quaisquer
$\vec a,\vec b$.

De maneira análoga, a espaçonave é descrita por um operador de inércia
espacial $M_{\text{espaç},\text{spt}}$.
No formalismo de velocidades espaciais \cite{featherstone2008}, uma matriz de
inércia espacial relaciona a velocidade espacial
$\vec v_{\text{sp}} = [\,{}^{B}\!\vec v^{\,T}\;{}^{B}\!\vec\omega^{\,T}\,]^{T}$
com a quantidade de movimento espacial
$\vec h_{\text{sp}}$ por meio de
$\vec h_{\text{sp}} = M_{\text{spt}}\,\vec v_{\text{sp}}$.

Para a espaçonave, $M_{\text{espaç},\text{spt}}$ é construído a partir da massa
$m_B$ e do matriz ${}^{B}\!I_B$ em torno do seu centro de massa, em completa
analogia com $M_{i,\text{spt}}$ definido para os elos.
A energia cinética total do conjunto espaçonave--\gls{mr} resulta da soma das
energias cinéticas de cada elo e da espaçonave.
Utilizando a forma espacial para cada elo, tem-se:

\begin{equation}
T
=
\sum_{i=1}^{n}
\frac{1}{2}
\bigl(J_{\text{sp},i}(\vec y\, )\,\dot{\vec y\, }\bigr)^{T}
M_{i,\text{spt}}
\bigl(J_{\text{sp},i}(\vec y\, )\,\dot{\vec y\, }\bigr)
+
\frac{1}{2}\,
\vec v_{\text{espaç,sp}}^{\,T}\,
M_{\text{espaç},\text{spt}}\,
\vec v_{\text{espaç,sp}}.
\end{equation}

Em que $\vec v_{\text{espaç,sp}}$ é a velocidade espacial da espaçonave, escrita
também como combinação linear das velocidades generalizadas $\dot{\vec y\, }$.
Comparando essa expressão com a forma quadrática de \eqref{eq_T_quadratica_Mq1}, obtém-se:

\begin{equation}
T =
\frac{1}{2}\,\dot{\vec y\, }^{\,T}\,M(\vec y\, )\,\dot{\vec y\, }.
\end{equation}

De modo que a matriz de inércia generalizada $M_s(\vec y\, )$ é dada pela soma das
contribuições espaciais de todos os elos, ponderadas pelos respectivos
jacobianos espaciais, acrescida da contribuição inercial da espaçonave:

\begin{equation}
M_s(\vec y\, ) =
\sum_{i=1}^{n}
J_{\text{sp},i}^{T}(\vec y\, )\,
M_{i,\text{spt}}\,
J_{\text{sp},i}(\vec y\, )
+
M_{\text{espaç},\text{spt}}.
\label{eq_Ms_total_formula}
\end{equation}

Cada termo
$J_{\text{sp},i}^{T}(\vec y\, \, \, )\,M_{i,\text{spt}}\,J_{\text{sp},i}(\vec y\, \, \, )$
é, portanto, a contribuição espacial do elo $i$ para a matriz de inércia
generalizada $M_s(\vec y\, )$.
Nessa expressão, cada jacobiano espacial $J_{\text{sp},i}(\vec y\, \, )$ mapeia as velocidades generalizadas $\dot{\vec y\, \, }$ na velocidade espacial (linear e angular) do centro de massa do elo $i$ em $\{B\}$.